\chapter{Green-Relationen und innere Struktur}
\label{b37guide:chap:green}

\section{Teilbarkeit und Stabilität}

\begin{definition}[Green-Relationen und Green-Präordnungen]
Für $a,b\in S$ seien
\[
\begin{aligned}
a\leq_{\mathcal R}b&\iff a\in bS^1,&
a\leq_{\mathcal L}b&\iff a\in S^1b,&
a\leq_{\mathcal J}b&\iff a\in S^1bS^1.
\end{aligned}
\]
Für $\mathcal K\in\{\mathcal R,\mathcal L,\mathcal J\}$ bedeutet
$a\mathrel{\mathcal K}b$, daß $a\leq_{\mathcal K}b$ und
$b\leq_{\mathcal K}a$ gelten. Ferner seien
$\mathcal H=\mathcal R\cap\mathcal L$ und $\mathcal D$ die kleinste
Äquivalenzrelation, die $\mathcal R\cup\mathcal L$ enthält.
Die Klassen von $a$ werden mit $R_a,L_a,H_a,D_a,J_a$ bezeichnet.
Auf den $\mathcal J$-Klassen gilt
$J_a\leq J_b\iff a\leq_{\mathcal J}b$.
\end{definition}

\begin{theorem}[Stabilität endlicher Halbgruppen]
\label{b37guide:thm:stability}
Es gilt $\mathcal J\subseteq\mathcal D$. Für $a\in S$ und $x\in S^1$ gilt
\[
a\mathrel{\mathcal J}ax\Longrightarrow a\leq_{\mathcal R}ax,
\qquad
a\mathrel{\mathcal J}xa\Longrightarrow a\leq_{\mathcal L}xa.
\]
\cite{pin,donoven2018}
\end{theorem}

\section{Struktur ganzer \texorpdfstring{$\mathcal J$}{J}-Klassen}

\begin{definition}[Reguläre Elemente und Untergruppen]
Ein Element $a\in S$ heißt \emph{regulär}, wenn ein $x\in S$ mit
$axa=a$ existiert. Eine $\mathcal J$-Klasse heißt \emph{regulär}, wenn
sie ein Idempotent enthält.

Eine \emph{Untergruppe von $S$} ist eine nichtleere Teilmenge $G\subseteq S$,
die produktabgeschlossen ist und ein Element $e\in G$ besitzt mit
$eg=ge=g$ für alle $g\in G$, wobei zu jedem $g\in G$ ein $h\in G$
mit $gh=hg=e$ existiert.
Ihr neutrales Element muß nur innerhalb von $G$ neutral sein.
Eine Untergruppe $G$ heißt \emph{zyklisch}, wenn es ein $g\in G$
mit $G=\{g^k:k\in\mathbb Z\}$ gibt; die Potenzen werden dabei
innerhalb der Gruppe gebildet.
Eine Untergruppe heißt \emph{maximal}, wenn sie in keiner größeren
Untergruppe von $S$ enthalten ist.
\end{definition}

\begin{theorem}[Klassenraster, Regularität und Gruppen innerhalb einer $\mathcal J$-Klasse]
\label{b37guide:thm:regularity}
In jeder $\mathcal J$-Klasse $J$ haben jede darin enthaltene
$\mathcal R$-Klasse und jede darin enthaltene $\mathcal L$-Klasse
einen nichtleeren Schnitt. Alle $\mathcal H$-Klassen in $J$ sind gleichmächtig.

Ist $J$ regulär, so ist jedes Element von $J$ regulär.
Je zwei maximale Untergruppen von $S$,
die in derselben $\mathcal J$-Klasse liegen, sind isomorph.
\cite{pin,izhakian2011}
\end{theorem}

\section{Lokale Halbgruppen an Idempotenten}

\begin{definition}[Lokale Halbgruppe]
Für $e\in E(S)$ heißt
\[
eSe=\{ese:s\in S\}
\]
mit der aus $S$ übernommenen Multiplikation die \emph{lokale Halbgruppe
bei $e$}. Die Bezeichnung bezieht sich auf das Idempotent $e$;
ein neutrales Element von $S$ wird nicht vorausgesetzt.
\end{definition}

\begin{theorem}[Einseitige Inversen und lokale Isomorphietypen]
\label{b37guide:thm:local-monoids}
Für jedes $e\in E(S)$ gilt
\[
(\forall x,y\in eSe)\quad xy=e\iff yx=e.
\]
Für $e,f\in E(S)$ gilt ferner
\[
e\mathrel{\mathcal J}f\Longrightarrow eSe\cong fSf.
\]
\cite{izhakian2011,howie}
\end{theorem}

\chapter{Ideale und Hauptfaktoren}
\label{b37guide:chap:ideals}

\section{Idealordnung und kleinstes Ideal}

\begin{definition}[Ideale und kleinstes Ideal]
Ein \emph{Linksideal} ist eine nichtleere Teilmenge $I\subseteq S$ mit
$SI\subseteq I$; ein \emph{Rechtsideal} erfüllt $IS\subseteq I$.
Ein \emph{Ideal} erfüllt beide Bedingungen. Ein Ideal $I$ heißt
\emph{kleinstes Ideal}, wenn $I\subseteq J$ für jedes Ideal $J$ von $S$ gilt.
\end{definition}

\begin{theorem}[Existenz und Struktur des kleinsten Ideals]
\label{b37guide:thm:kernel}
Für $S\ne\varnothing$ besitzt $S$ ein kleinstes Ideal $K(S)$.
Jede $\mathcal H$-Klasse in $K(S)$ ist eine Gruppe.
Für jedes $e\in K(S)\cap E(S)$ gilt $eSe\subseteq H_e$.
\cite{east2018}
\end{theorem}

\section{Lokale Multiplikationsstruktur}

\begin{definition}[Hauptfaktor einer $\mathcal J$-Klasse]
Für eine $\mathcal J$-Klasse $J$ ist ihr \emph{Hauptfaktor}
$P(J)=J$, wenn $J=K(S)$, mit der Multiplikation aus $S$.
Andernfalls ist $P(J)=J\cup\{0_J\}$, wobei $0_J$ neu ist und
\[
x\circ y=
\begin{cases}
xy,&x,y\in J\text{ und }xy\in J,\\
0_J,&\text{sonst}.
\end{cases}
\]
Eine \emph{Nullhalbgruppe} ist eine Halbgruppe mit einem Element $0$,
in der jedes Produkt gleich $0$ ist.
\end{definition}

\clearpage
\begin{definition}[Rees-Matrixhalbgruppen]
Seien $G$ eine endliche Gruppe, $A,B$ nichtleere endliche Mengen und
$P=(p_{\beta a})_{\beta\in B,a\in A}$ eine Matrix über
$G\cup\{0\}$, wobei $0\notin G$ ein formales Zeichen ist.
Die \emph{Rees-Nullmatrixhalbgruppe}
$M^0(G;A,B;P)=(A\times G\times B)\cup\{0\}$ hat $0$ als
absorbierendes Element und die Multiplikation
\[
(a,g,\beta)(c,h,\delta)=
\begin{cases}
(a,gp_{\beta c}h,\delta),&p_{\beta c}\ne0,\\
0,&p_{\beta c}=0.
\end{cases}
\]
Die Matrix heißt \emph{regulär}, wenn jede Zeile und jede Spalte
einen von $0$ verschiedenen Eintrag enthält.
Hat $P$ nur Einträge aus $G$, bezeichnet $M(G;A,B;P)$ die
Unterhalbgruppe $A\times G\times B$ mit derselben Produktformel.
\end{definition}

\begin{theorem}[Rees--Suschkewitsch-Struktur der Hauptfaktoren]
\label{b37guide:thm:rees}
Für jede $\mathcal J$-Klasse $J$ gilt einer der folgenden Fälle:
\begin{enumerate}
\item $J$ ist nichtregulär. Dann ist $J\ne K(S)$ und $P(J)$ eine
Nullhalbgruppe.
\item $J$ ist regulär und $J\ne K(S)$. Dann ist
\[
P(J)\cong M^0(G;A,B;P)
\]
für geeignete nichtleere endliche Mengen $A,B$, eine endliche Gruppe
$G$ und eine reguläre Matrix $P$.
\item $J=K(S)$. Dann ist
\[
K(S)\cong M(G;A,B;P)
\]
für geeignete nichtleere endliche Mengen $A,B$, eine endliche Gruppe
$G$ und eine Matrix $P$ mit ausschließlich Einträgen aus $G$.
\end{enumerate}
\cite{rees1940,shah2016}
\end{theorem}

\clearpage
\section{Maximale echte Unterhalbgruppen}

\begin{definition}[Maximalität und Rees-Koordinaten]
Eine nichtleere, produktabgeschlossene Teilmenge $M\subsetneq S$ heißt
\emph{maximale echte Unterhalbgruppe}, wenn keine Unterhalbgruppe $N$
mit $M\subsetneq N\subsetneq S$ existiert. Entsprechend heißt eine
Untergruppe $H\subsetneq G$ einer Gruppe $G$ \emph{maximal echt}, wenn
keine Untergruppe $K$ mit $H\subsetneq K\subsetneq G$ existiert.

Für eine reguläre $\mathcal J$-Klasse $J$ heißen die durch Einschränkung
eines Rees-Isomorphismus aus Satz~\ref{b37guide:thm:rees} entstehenden
Bijektionen
\[
c:J\longrightarrow A\times G\times B
\]
\emph{Rees-Koordinatisierungen} von $J$.
\end{definition}

\begin{theorem}[Form maximaler Unterhalbgruppen nach Graham--Graham--Rhodes]
\label{b37guide:thm:maximal-subsemigroups}
Sei $M$ eine maximale echte Unterhalbgruppe der endlichen Halbgruppe $S$.
Für ein beliebiges $a\in S\setminus M$ setzen wir $J=J_a$.
Dann tritt einer der folgenden Fälle ein; in den regulären Fällen
können geeignete, von $M$ abhängige Rees-Koordinaten
$c:J\to A\times G\times B$ gewählt werden:
\begin{enumerate}
\item $J$ ist nichtregulär und $M\cap J=\varnothing$.
\item $J$ ist regulär, und für nichtleere Teilmengen $A'\subseteq A$
und $B'\subseteq B$ gilt
\[
c(M\cap J)=(A\times G\times B)\setminus(A'\times G\times B').
\]
\item $J$ ist regulär, und für eine maximale echte Untergruppe $H$ von
$G$ gilt
\[
c(M\cap J)=A\times H\times B.
\]
\end{enumerate}
\cite{graham1968,donoven2018}
\end{theorem}

\section{Endliche Gliederung durch Hauptfaktoren}

\begin{definition}[Hauptreihe]
Eine \emph{Hauptreihe} ist eine Kette
$\varnothing=I_0\subsetneq I_1\subsetneq\cdots\subsetneq I_m=S$,
in der $I_1,\ldots,I_m$ Ideale sind und zwischen zwei aufeinanderfolgenden
Gliedern kein weiteres Ideal liegt. Für $S=\varnothing$ besteht die
Hauptreihe nur aus $I_0=\varnothing$.
\end{definition}

\begin{theorem}[Endliche Gliederung durch Hauptfaktoren]
\label{b37guide:thm:principal-series}
$S$ besitzt eine endliche Hauptreihe.
\cite{howie}
\end{theorem}

\section{Überdeckung durch zwei Unterhalbgruppen}

\begin{definition}[Monogene Halbgruppe]
Eine Halbgruppe $A$ heißt \emph{monogen}, wenn es ein $a\in A$ mit
\[
 A=\langle a\rangle=\{a^k:k\geq1\}
\]
gibt.
\end{definition}

\begin{theorem}[Überdeckungssatz von Donoven--Kappe]
\label{thm:b37-guide-two-subsemigroup-cover}
Für jede nichtleere endliche Halbgruppe $S$ gilt mindestens eine
der folgenden Alternativen:
\begin{enumerate}
 \item $S$ ist monogen.
 \item $S$ ist eine Gruppe.
 \item Es gibt echte Unterhalbgruppen $A,B\subsetneq S$ mit
 \[
 S=A\cup B.
 \]
\end{enumerate}
\cite[Satz~1.4]{DonovenKappe2023}
\end{theorem}
